\newcommand{\cvexperience}{%
\cventry{feb.--ago. 2026}
  {Microservicios Reactivos con WebFlux y Kafka}
  {Smart Financial SAS}{Cali, Colombia}{
  \begin{itemize}
    \item Implement\'e microservicios bajo arquitectura reactiva, contribuyendo a que el sistema pudiera gestionar picos de carga significativos sin afectar su estabilidad.
    \item Apliqu\'e persistencia no bloqueante para evitar cuellos de botella de E/S, lo que represent\'o un cambio significativo en la capacidad de respuesta del sistema.
    \item Desacoplé procesos críticos mediante mensajer\'{\i}a as\'incrona, logrando que tareas en segundo plano no interrumpieran el flujo transaccional principal y generando un cambio significativo en la fluidez del procesamiento.
    \item Contribuí a la configuraci\'on y gobernanza del API Gateway, reforzando el control de acceso y la protecci\'on ante solicitudes excesivas, lo que supuso un cambio significativo en la gesti\'on de los servicios.
    \item Integr\'e métricas y monitoreo en tiempo real, facilitando la detecci\'on temprana de anomal\'{\i}as y contribuyendo a un aumento significativo en la estabilidad operativa.
    \item Realic\'e pruebas de carga escalonada para validar el comportamiento del sistema bajo escenarios exigentes, lo que permiti\'o identificar ajustes y generar cambios significativos en su capacidad para soportar alta concurrencia.
    \item Apoyé la documentaci\'on arquitect\'onica y el an\'alisis de rendimiento del sistema, lo que facilit\'o decisiones t\'ecnicas que representaron cambios significativos en la mantenibilidad y robustez del proyecto.
    \item Nota: Kubernetes y Kong API Gateway fueron implementados por el equipo de DevOps; el resto lo desarroll\'e directamente.
  \end{itemize}
}{}
}
